\documentclass[11pt,a4paper]{article}
\usepackage[utf8]{inputenc}
\usepackage[T1]{fontenc}
\usepackage{geometry}
\geometry{margin=2.5cm}
\usepackage{booktabs}
\usepackage{hyperref}
\usepackage{enumitem}

\title{Prompts Log: Knowledge Engineering Tasks}
\author{Lucas Perez Reis Lobo, Yutaka Shibata, Sofia Davis, \\ William Caulton, Mert Yerlikaya, Davyd Shtepa \\ \small Knowledge Engineering CW2 --- Music History \& Heritage KG}
\date{April 2026}

\begin{document}
\maketitle

\section{Overview}

43 prompts were used across 7 knowledge engineering tasks. All used Claude (Opus 4.6) or Google Gemini (free tier).

KE Tasks: Ontology and data source discovery (P1, P2, P7, 3 prompts), Competency question generation (P3--P6, 4 prompts), LLM text extraction mapping (P8--P13, P24--P33, 20 prompts), RAG gap-filling (P14--P17, 4 prompts), KG vs Plain LLM evaluation (P18--P23, 6 prompts), RAG completion of 10 gaps (P34--P43, 10 prompts). Total: 43 prompts.

\section{Ontology and Data Source Discovery (P1, P2, P7)}

\textbf{P1:} ``Research existing ontologies related to music history and arts/cultural heritage. I need at least 2 that can be extended with subclasses and subproperties.''

\textbf{P2:} ``Research data sources for a music history KG. I need at least 1 textual and 1 structured source. Search for: MusicBrainz, Spotify, Wikidata, Discogs, Last.fm, Wikipedia, BBC Music.''

\textbf{Justification:} Open-ended prompts provided breadth at this early stage. The constraint in P1 (``can be extended with subclasses and subproperties'') ensured results focused on extensible ontologies rather than closed vocabularies.

\textbf{Output use:} P1 identified 4 candidates (Music Ontology, CIDOC-CRM, DOREMUS, Schema.org) with concrete extension suggestions; the team selected Music Ontology and Schema.org. P2 identified 5 structured and 3 textual sources; we selected MusicBrainz, Discogs, and Wikipedia. Rejected: Spotify (requires OAuth), Last.fm (registration blocked).

\textbf{Strategy refinement:} After 7 experiments, P7 re-evaluated with a targeted prompt: ``Are these the best data sources? Test across genres: classical, jazz, rock, Brazilian, African.'' This led to adding Wikidata SPARQL (awards and instruments already in RDF, eliminating the need for LLM extraction of these properties) and rejecting DBpedia (template parameters leaking as properties) and Open Opus (returns ``NOT FOUND'' for non-classical artists).

\section{Competency Question Generation (P3--P6)}

\textbf{P5 (key prompt):} ``You are a knowledge engineer and music historian. Generate 10 CQs that complement these 10 manual CQs: [list]. Cover: record labels, collaborations, influence chains, geographic spread. At least 3 require 3+ entity joins. Include relationship path for each. [3 examples with paths].''

\textbf{Justification:} Combined role-based (domain grounding), few-shot (3 examples set complexity), and constrained output (explicit requirements).

\textbf{Output use:} P5 generated 10 CQs, all implemented as SPARQL queries. 8/10 require 3+ entity joins. Covered gaps: record labels (CQ11,13,17), collaborations (CQ14,19), influences (CQ12,18), geography (CQ16,20).

\textbf{Strategy refinement:} The few-shot examples in P5 caused the LLM to copy rather than generate new questions. P6 fixed this: ``Review the 10 CQs. Think step by step. Issues: (1)~CQs 1--3 copy my few-shot examples. (2)~Vague aggregations are hard for SPARQL. (3)~None test the MusicalWork to Recording distinction. (4)~Verify each is unique and SPARQL-implementable.'' P6 replaced 3 weak CQs while preserving 7 good ones --- more efficient than regenerating all 10. The copying issue informed later prompts: Batch~3 used targeted emphasis instead of few-shot examples.

\section{LLM Text Extraction (P8--P33, 20 prompts)}

\textbf{Template:} ``You are a knowledge engineer and music historian. Extract structured triples from this Wikipedia article about [ARTIST]. Use ONLY: released, composed, collaboratedWith, alterEgo, albumGrouping, producedBy, pioneerOf, founded, hasGenre, performedAt. Do NOT extract: birth/death dates, awards, instruments, labels. Return as JSON array.''

\textbf{Justification:} Controlled predicates ensure output maps to our ontology. The ``Do NOT extract'' list prevents duplication with Wikidata. JSON enables programmatic ingestion.

\textbf{Output use} (P8 David Bowie example): [\{``subject'':``David Bowie'',``predicate'':``alterEgo'',``object'':``Ziggy Stardust''\}, \{``subject'':``David Bowie'',``predicate'':``hasGenre'',``object'':``plastic soul''\}, \{``subject'':``Heroes (1977)'',``predicate'':``albumGrouping'',``object'':``Berlin Trilogy''\}]. These facts exist only in Wikipedia text. 897 triples extracted across 58 artists. 44\% use predicates no structured API provides.

\textbf{Strategy refinement:} Batch~1 (P8--P13, 103 triples) used 8 predicates. Analysis showed sparse results for producedBy, pioneerOf, performedAt. Batch~2 (P24--P28, 93 triples) added these predicates. Batch~3 (P29--P33, 102 triples) added ``Focus ESPECIALLY on: producedBy, albumGrouping, founded'' --- increasing producedBy from 1 to 15 triples. The template was then automated via Google Gemini API, reaching 58/58 coverage.

\section{RAG Gap-Filling (P14--P17)}

\textbf{P14:} ``Our KG contains Michael Jackson's discography but is missing producer information. Context: [8 albums with dates from SPARQL]. Our KG already knows Quincy Jones produced Off the Wall, Thriller, and Bad. For each album, who was the primary producer? Return as JSON with confidence levels.''

\textbf{Justification:} RAG grounds responses in actual KG data. A plain prompt would generate answers from parametric knowledge, potentially including album titles not in our KG.

\textbf{Output use:} The LLM confirmed 3 existing triples, added 2 high-confidence new ones, and flagged 3 as uncertain (multiple co-producers). Only high-confidence triples were ingested. P16 resolved 14/15 missing countries. P17 used influence-chain context to assign genres to influencers.

\textbf{Strategy refinement:} The LLM's unprompted confidence self-assessment in P14 was unexpected. We adopted this pattern for all subsequent RAG prompts (P34--P43), always requesting confidence levels.

\section{KG vs Plain LLM Evaluation (P18--P23)}

Six paired comparisons testing the same question as (1)~plain zero-shot and (2)~RAG with SPARQL-retrieved KG context (e.g., ``Our KG shows 34 awards: [list]. Are these complete and accurate?''). P18 RAG response: ``Critical warning for KG integrity: the number 34 is plausible but not verifiable. Inserting this value as true without cross-referencing would introduce untracked uncertainty.'' The LLM refused to fabricate missing awards. P18 informed award blacklist. P20 led to instrument normalisation. P21 led to genre blacklist additions. P22 blacklisted ``classic rock''. Every comparison produced at least one actionable fix. RAG functioned as a data quality auditor.

\section{RAG Completion (P34--P43)}

10 prompts resolving the 5+5 gaps from the completion analysis. Each followed the 6-step RAG cycle. P35 example: SPARQL retrieved 58 primary artists and 5 existing performedAt triples. Prompt: ``For each of these 58 artists, what are their 2--3 most iconic live performances? Only include historically documented events. Return as JSON.'' The LLM appropriately skipped classical composers, producers, and studio artists. Total across all 10 prompts: 494 triples ingested. performedAt 5 to 131, collaboratedWith 76 to 242, founded 4 to 69, hasMusicalPeriod 5 to 126. P40 showed the LLM hallucinating a MusicBrainz ID (returned 404). P42 compensated for structured source blind spots in medieval music.

\section{Strategy Evolution}

The prompting approach matured through four phases. Early prompts (P1--P2) used open-ended research queries suitable for discovery. The CQ phase (P3--P6) introduced structured techniques, but few-shot examples backfired by causing copying, teaching us that examples constrain creativity. This led to targeted emphasis in Batch~3, which guided attention without constraining output. The RAG phase (P14--P43) was the most significant shift: grounding prompts in SPARQL-retrieved KG data transformed the LLM from a content generator into a data quality auditor. The discovery that RAG reveals KG errors (not just fills gaps) reshaped our approach, leading to 6 pipeline fixes and the automation of text extraction via API.

Full prompt details: Complete prompt text and model responses for all 43 prompts are available in the repository at \texttt{docs/prompts\_log.md} (\url{https://github.com/LucasPRLobo/ke-cw2}).

\textbf{Validation:} All LLM outputs were verified before ingestion into the KG. Text-extracted triples passed through a 5-tier entity linking system achieving 99\% accuracy. RAG-ingested triples were validated by 22 structural rules. The final KG was reviewed in Prot\'{e}g\'{e} for ontological correctness and by a domain expert for semantic accuracy. External identifiers from LLM responses were always verified programmatically, after P40 demonstrated that LLMs hallucinate plausible-looking but nonexistent IDs.

\end{document}
